% Generated draft: requires mathematical and source review.
\documentclass[11pt]{article}
\usepackage[margin=25mm]{geometry}
\usepackage{fontspec}
\setmainfont{Latin Modern Roman}
\usepackage{amsmath,amssymb,mathtools,unicode-math}
\setmathfont{Latin Modern Math}
\usepackage{enumitem,needspace,xurl,hyperref}
\setlength{\parindent}{0pt}
\setlength{\parskip}{5pt}
\setlength{\emergencystretch}{3em}
\newcommand{\N}{\mathbb{N}}
\newcommand{\Z}{\mathbb{Z}}
\newcommand{\Q}{\mathbb{Q}}
\newcommand{\R}{\mathbb{R}}
\newcommand{\C}{\mathbb{C}}
\newcommand{\F}{\mathbb{F}}
\newcommand{\A}{\mathbb{A}}
\newcommand{\PP}{\mathbb{P}}
\newcommand{\E}{\mathbb{E}}
\newcommand{\eps}{\varepsilon}
\newcommand{\dd}{\,\mathrm{d}}
\begin{document}
\section*{UnsolvedMath ID: 2966}
\textbf{Source problem code:} catalogue entry
\subsection{Definitions and mathematical statement}
Let \((X,\omega)\) be a closed symplectic 4-manifold with its symplectic orientation. For an \(\omega\)-compatible almost complex structure \(J\), define
\[
c_1^2(X)=\langle c_1(TX,J)^2,[X]\rangle,\qquad
c_2(X)=\langle c_2(TX,J),[X]\rangle.
\]
The supplied identities are \(c_1^2=2\chi(X)+3\sigma(X)\) and \(c_2=\chi(X)\), where \(\sigma\) is the signature with that orientation. The choice of compatible \(J\) does not change these numbers.

For every such \(X\) outside the intended ruled-surface class, the three source questions are:
\begin{enumerate}
\item \(c_1^2(X)\leq3c_2(X)\).
\item \(c_1^2(X)=3c_2(X)>0\) if and only if \(X\) is \(\mathbb{CP}^2\) or a complex ball quotient.
\item \(c_2(X)\geq0\).
\end{enumerate}
A closed smooth complex ball quotient means \(\mathbb B^2_\C/\Gamma\), where \(\Gamma\) is a torsion-free cocompact discrete group of holomorphic isometries of the complex unit ball. Since \(X\) initially has only a symplectic structure, the meaning of ``is'' in part (b) must be fixed, for example oriented diffeomorphism type rather than equality of chosen complex structures. The source does not state this equivalence explicitly. Connectedness is also implicit in the usual equality classification; if disconnected manifolds are allowed, a componentwise formulation is needed.

The meaning of ruled must likewise be fixed: an \(S^2\)-bundle over a closed surface, versus a class also including its blowups, yields different domains. The background uses K\"ahler ruled-surface language and a separate minimality discussion but does not give an explicit definition. No extra minimality or \(b_2^+>1\) restriction is silently added to the main three questions. Minimal in that discussion means that there is no symplectically embedded sphere of self-intersection \(-1\).

The source reports the three claims for K\"ahler surfaces, with part (a) the Bogomolov--Miyaoka--Yau inequality and the equality case linked to Yau's theorem. It attributes part (c)'s symplectic conjecture to Gompf and reports \(c_1^2\geq0\) for minimal nonruled symplectic manifolds from Taubes and Li--Liu, so (a) would imply (c) in that case. It describes potential consequences for exotic symplectic manifolds homeomorphic to \(\mathbb{CP}^2\) or to ruled surfaces over genus at least two, without claiming those consequences established. It also asks whether hypothetical non-K\"ahler equality examples might be aspherical, as a separate direction associated with Problem 4.91.

The background reports Auroux's branched-cover approach, expressing geography through a singular symplectic branch curve in \(\mathbb{CP}^2\), with nodes, simple cusps and covering data; no missing numerical formula is invented. It proposes the broader Seiberg--Witten question for closed manifolds with \(b_2^+>1\) and nontrivial invariants, and describes Feehan--Leness's announced program using non-Abelian monopoles and adaptations of Hitchin's Morse-theoretic analysis. These are approaches and proposed extensions, not completed proofs of (a).

The registry status is open. The supplied triage dated 2026-08-17 labels the row partially solved, records the topological identities and complex ball equality examples, and reports no full symplectic proof of the sharp inequalities or equality classification. No announced program, K\"ahler theorem or current status is independently verified here.
\subsection{Short English statement}
For closed nonruled symplectic 4-manifolds, ask for the BMY Chern-number inequality, its positive equality classification, and nonnegative Euler characteristic. The source records K\"ahler results and symplectic partial constraints, but the ruled and equality-equivalence conventions must be explicit, and announced proof programs are not full solutions.
\subsection{Sources}
\begin{itemize}
\item Local UnsolvedMath registry, permanent ID 2966, source code catalogue entry, ``Kirby Problem 4.90,'' including supplied literature triage dated 2026-08-17 and related Problem 4.91 and 4.16 pointers.
\item \emph{K3: A New Problem List in Low-Dimensional Topology}, Problem 4.90 (2026): \url{https://aimath.org/pastworkshops/kirbylistrep.pdf}.
\item W. Barth, C. Peters, and A. Van de Ven, \emph{Compact Complex Surfaces}, Ergebnisse der Mathematik (3), vol. 4, Springer-Verlag, Berlin (1984), DOI 10.1007/978-3-642-96754-2, source label [BPVdV84].
\item Shing Tung Yau, ``On the Ricci curvature of a compact K\"ahler manifold and the complex Monge-Amp\`ere equation. I,'' \emph{Communications on Pure and Applied Mathematics} 31(3) (1978), 339--411, DOI 10.1002/cpa.3160310304, source label [Yau78].
\item Clifford Henry Taubes, ``The Seiberg-Witten invariants and symplectic forms,'' \emph{Mathematical Research Letters} 1(6) (1994), 809--822, DOI 10.4310/MRL.1994.v1.n6.a15, source label [Tau94].
\item Clifford Henry Taubes, ``More constraints on symplectic forms from Seiberg-Witten invariants,'' \emph{Mathematical Research Letters} 2(1) (1995), 9--13, DOI 10.4310/MRL.1995.v2.n1.a2, source label [Tau95].
\item Clifford H. Taubes, ``\(\mathrm{SW}\Rightarrow\mathrm{Gr}\): from the Seiberg-Witten equations to pseudo-holomorphic curves,'' \emph{Journal of the American Mathematical Society} 9(3) (1996), 845--918, DOI 10.1090/S0894-0347-96-00211-1, source label [Tau96]. The source title's arrow glyph is damaged.
\item T. J. Li and A. Liu, ``Symplectic structure on ruled surfaces and a generalized adjunction formula,'' \emph{Mathematical Research Letters} 2(4) (1995), 453--471, DOI 10.4310/MRL.1995.v2.n4.a6, source label [LL95].
\item Denis Auroux, ``Symplectic 4-manifolds as branched coverings of \(\mathbb{CP}^2\),'' \emph{Inventiones Mathematicae} 139(3) (2000), 551--602, DOI 10.1007/s002220050019, source label [Aur00].
\item Denis Auroux, ``Symplectic 4-manifolds, singular plane curves, and isotopy problems,'' in \emph{Floer Homology, Gauge Theory, and Low-Dimensional Topology}, Clay Mathematics Proceedings 5, American Mathematical Society (2006), 263--276, supplied pointer p. 266, source label [Aur06b].
\item Paul M. N. Feehan, ``Bialynicki-Birula theory, Morse-Bott theory, and resolution of singularities for analytic spaces,'' arXiv:2206.14710 (2022), source label [Fee22].
\item Paul M. N. Feehan and Thomas G. Leness, ``Virtual Morse-Bott index, moduli spaces of pairs, and applications to topology of smooth four-manifolds,'' arXiv:2010.15789, dated 2023 in the supplied bibliography, source label [FL23].
\item Paul M. N. Feehan and Thomas G. Leness, ``Almost Hermitian structures on moduli spaces of non-abelian monopoles and applications to the topology of symplectic four-manifolds,'' arXiv:2410.13809 (2024), source label [FL24b].
\item N. J. Hitchin, ``The self-duality equations on a Riemann surface,'' \emph{Proceedings of the London Mathematical Society} (3) 55(1) (1987), 59--126, DOI 10.1112/plms/s3-55.1.59, source label [Hit87].
\end{itemize}
\end{document}
